 \documentclass[a4paper,10pt]{article}
\usepackage{paper-ru}
\usepackage{hyperref}

\def\thetitle{Дидона с велосипедом}
\def\theauthors{Антон Петрунин}

\hypersetup{colorlinks=true,
citecolor=black,
linkcolor=black,
anchorcolor=black,
filecolor=black,
menucolor=black,
urlcolor=black,
pdftitle={\thetitle},
pdfauthor={\theauthors}
}

%\overfullrule=100mm


\begin{document}
\title{\thetitle}
\author{\theauthors}
\date{}
\maketitle

Мы покажем, как, используя велосипед, можно решить задачу Дидоны.
То есть доказать, что из всех фигур, зажатых между кривой длины не больше~$L$ и прямой через её концы полукруг имеет наибольшую площадь.

Это первая известная изопериметрическая задача.
Царице Дидоне пришлось отгораживать ремнём из бычьей шкуры участок у берега моря под город Карфаген, и это случилось в конце девятого века до нашей эры.
Обычная изопериметрическая задача говорит, что из всех замкнутых кривых на плоскости данной длины окружность ограничивает наибольшую площадь.
Эта задача легко сводится к задаче Дидоны и наоборот --- попробуйте в этом убедиться сами.

\begin{figure}[ht!]
\vskip-0mm
\centering
\includegraphics{mppics/pic-1}
\end{figure}

Нам понадобится математический велосипед.
Будем считать, что рама велосипеда --- это отрезок постоянной длины (скажем, $R$); задний его конец ($Z$) может двигаться только в направлении рамы, а передний ($P$) в любом направлении.
Это вполне согласуется с практикой: если исключить проскальзывание, то заднее колесо обычного велосипеда движется вдоль рамы, а переднее можно направить как угодно.


Предположим, что переднее колесо движется с единичной скоростью.
В этом случае составляющая его скорости, перпендикулярная раме, не превосходит $1$.
Поскольку заднее колесо движется вдоль рамы, угловая скорость рамы не может превышать $1/R$ по модулю.

\begin{wrapfigure}{r}{30 mm}
\vskip-0mm
\centering
\includegraphics{mppics/pic-3}
%\caption*{}
\end{wrapfigure}

Теперь представим, что мы закрашиваем область под рамой велосипеда;
то есть закрашивается объединение всех положений рамы.
За короткий промежуток времени $\Delta t$ (в первом приближении) закрасится треугольник с основанием $R$ и высотой не больше $\Delta t$,
а значит скорость закраски не превысит $R/2$.

Предположим, что переднее колесо прошло вдоль кривой длины~$L$, и при этом закрасилась площадь $S$, а рама повернулась на угол $\Theta$; угол измеряется в радианах.
Из вышесказанного получаем, что
\[\Theta\le L/R\qquad\text{и}\qquad S\le R\cdot L/2.\leqno({*})\]
Эти наблюдения сыграют ключевую роль в нашем доказательстве.

\medskip

Пусть $F$ --- фигура, зажатая между кривой $\gamma$ длины не больше $L$ с концами на прямой и самой прямой.
Далее пусть $D$ --- полукруг, длина дуги которого равна $L$; значит, его радиус равен $R=L/\pi$, а площадь $R\cdot L/2=\pi\cdot R^2/2$.
Нам надо доказать, что площадь $F$ не превосходит площади~$D$.

\begin{wrapfigure}[5]{r}{40 mm}
\vskip-4mm
\centering
\includegraphics{mppics/pic-5}
%\caption*{}
\end{wrapfigure}

Предположим, что это неверно, и фигура $F$ --- контрпример.
Фигуру $F$ можно считать выпуклой: иначе, выпрямляя вогнутые участки $\gamma$, как показано на рисунке, можно увеличить площадь и одновременно уменьшить длину, так что в результате мы получим выпуклый контрпример.

Далее можно считать, что граничная кривая $\gamma$ гладкая.
Это позволит нам проехаться вдоль $\gamma$ передним колесом велосипеда с единичной скоростью.
Действительно, если слегка уменьшить фигуру гомотетией, её площадь останется больше площади $D$, а длина кривой станет строго меньше $L$.
После этого кривую можно сгладить достаточно малым возмущением, оставив её выпуклой.

Теперь рассмотрим следующие две картинки.
На левой изображён полукруг $D$.
Представим, что мы провели переднее колесо велосипеда вдоль его дуги и в начальный момент рама образовывала радиус полукруга.
В этом случае достигаются оба равенства в $({*})$.
Действительно, рама будет оставаться радиусом во время всего движения, а заднее колесо останется в центре.
Поэтому рама повернётся на угол \(\pi=L/R\), и, проехав вдоль дуги, мы закрасим весь полукруг, площадь которого равна $R\cdot L/2=\pi\cdot R^2/2$.


\begin{figure}[h!]
\includegraphics{mppics/pic-10}
\end{figure}

На правой картинке мы проделали то же для кривой $\gamma$.
Опять же, считаем, что в начальный момент рама шла вдоль граничного отрезка так, что вся фигура лежала справа от рамы.
Если мы убедимся, что вся фигура $F$ закрасится после проезда передним колесом вдоль $\gamma$, второе неравенство в $({*})$ приведёт к противоречию и завершит доказательство.

\medskip

Предположим, что во время всего пути руль был направлен прямо или вправо, а значит рама велосипеда не могла поворачиваться против часовой стрелки и была направлена наружу из фигуры.
По первому неравенству в $({*})$ за время всего пути рама повернулась не больше чем на $\pi=L/R$.
То есть в последний момент она располагалась на прямой или ниже.
В самом конце руль также не мог быть повёрнут влево, а значит вся фигура $F$ лежала слева от рамы.

Выберем точку $x$ в $F$.
В начале движения точка $x$ лежит справа от рамы, а в конце слева от неё.
Поэтому в некий промежуточный момент точка переходит с правой стороны рамы на левую.
В этот момент точка должна лежать на прямой, проходящей вдоль рамы,
и она может лежать на самой раме либо на одном из двух лучей, продолжающих раму: позади заднего колеса или впереди переднего.
В первом случае точка $x$ закрашена.
Второй случай невозможен, ведь рама не поворачивалась влево, а значит переход происходил бы не справа налево, а наоборот.
Последний случай также невозможен, ведь луч перед рамой всё время указывает наружу фигуры $F$, а поскольку фигура выпуклая, весь луч лежит снаружи.

Итак, любая точка $x$ фигуры закрашена, а значит, закрашена вся фигура.

Осталось разобрать случай, когда руль меняет направление поворота.
То есть в какой-то момент он был повёрнут влево, и значит рама направлена внутрь $F$, как показано на рисунке.
Такое действительно может случиться: на последней картинке показан такой вариант развития событий.

{

\begin{wrapfigure}{r}{34 mm}
\vskip-2mm
\centering
\includegraphics{mppics/pic-15}

\bigskip

\includegraphics{mppics/pic-20}
%\caption*{}
\end{wrapfigure}

Как и раньше, мы считаем, что фигура располагается справа по ходу кривой $\gamma$.
Посмотрим на то, как меняется угол $\alpha$ между рамой и направлением движения.

Пока переднее колесо повёрнуто влево относительно рамы, рама вращается против часовой стрелки.
В то же время направление касательной к выпуклой кривой $\gamma$ вращается по часовой стрелке.
Поэтому угол $\alpha$ не может уменьшаться.
Значит, в какой-то момент угол $\alpha$ оказался нулевым.
В этот момент переднее колесо движется точно в сторону заднего,
то есть велосипед едет назад.

Пусть $t_0$ --- первый такой момент времени.
Тогда вся фигура закрасилась уже к моменту $t_0$.
Действительно, в начале вся фигура лежит справа от рамы, а в момент времени $t_0$ --- слева.
Между этими моментами руль был повёрнут вправо или точно прямо --- осталось повторить доказательство основной части.

}

\paragraph{Замечания.}
Идея этого доказательства принадлежит Гэри Лолору~\cite{lawlor};
она применима и в случае равенства, хотя требуется дополнительная работа.
Попробуйте это проделать, то есть доказать, что если потребовать дополнительно, чтобы площади $F$ и $D$ были равны, то $F$ конгруентна $D$.


Попытайтесь разобраться, как и когда по следам велосипеда можно определить направление его движения, а также различить следы переднего и заднего колеса.
Это заставит вас переоткрыть ощутимую часть дифференциальной геометрии кривых, а кроме того, вы получите удовольствие.

Математический велосипед активно изучается, и с ним связано большое число открытых вопросов~\cite{foote-levi-tabachnikov,levi-tabachnikov}.
Например, предположим, что след переднего колеса велосипеда совпадает со следом заднего и при этом велосипед не едет по прямой.
Обязан ли такой след иметь самопересечения?


{\sloppy\printbibliography[heading=bibintoc]\fussy}
\end{document}

